% ATTEMPT_STATUS: unresolved
% TOP_PROBLEM: {"problemId": "problem.uniqueness-of-tangent-cones-for-area-minimizing-integral-currents-in-higher-codimension", "canonicalTitle": "Uniqueness of Tangent Cones for Area-Minimizing Integral Currents in Higher Codimension", "releaseRank": 418, "primaryDomain": "analysis_pde", "primaryDomainLabel": "Analysis & PDE", "displayStatus": "open_with_solved_subcases", "releaseStatus": "open_with_solved_subcases"}
% !TeX program = lualatex
\documentclass[11pt]{article}
\usepackage[margin=25mm]{geometry}
\usepackage{fontspec}
\setmainfont{Latin Modern Roman}
\usepackage{amsmath,amssymb,mathtools}
\usepackage{unicode-math}
\setmathfont{Latin Modern Math}
\usepackage{xurl,hyperref}
\hypersetup{colorlinks=true,urlcolor=blue}
\setcounter{secnumdepth}{0}
\setlength{\parindent}{0pt}
\setlength{\parskip}{5pt}
\setlength{\emergencystretch}{3em}
\begin{document}
\section*{\texorpdfstring{Uniqueness of Tangent Cones for Area-Minimizing Integral Currents in Higher Codimension}{Uniqueness of Tangent Cones for Area-Minimizing Integral Currents in Higher Codimension}}\label{418}
\subsection{Definitions and mathematical statement}
An integral $m$-current represents an oriented rectifiable set with integer multiplicity and locally finite mass, with integral boundary. It is area-minimizing if no compactly supported competitor with the same boundary has smaller mass. For an interior point $p\in\operatorname{spt}T\setminus\operatorname{spt}\partial T$ and $r>0$, define
\[
 T_{p,r}=(\eta_{p,r})_\#T,
 \qquad\eta_{p,r}(x)=\frac{x-p}{r}.
\]
A tangent cone is a local weak limit along some $r_j\downarrow0$; weak convergence tests against compactly supported smooth differential forms. For $m\ge3$ and codimension $n\ge2$, the conjecture asserts that every such subsequential limit is the same current at every interior point. It is uniqueness of the full oriented multiplicity-bearing tangent current, not merely of its support or density, and concerns minimizing currents rather than arbitrary stationary varifolds.
\par\smallskip\textit{Formulation and concept references:} \href{https://proceedings.centre-mersenne.org/item/10.5802/slsedp.175.pdf}{[S1]}.

\subsection{Short English statement}
At every interior point of an area-minimizing integral current in higher codimension, do all infinitesimal blow-up sequences converge to the same tangent cone?
\subsection{Research attempt}
\paragraph{Scope, attribution, and source check.}
This unresolved attempt was prepared by GPT-6 Astra Ultra on 27 September
2026 using the primary account [S1], explicit blow-up calculations and
checks of the quantifiers. The target is uniqueness of the oriented
integral current, including multiplicity, at every interior point.
The source describes uniqueness and decay results on specified strata
and recalls the two-dimensional theorem. These do not establish the
assertion at every point in dimensions $m\geq3$. The attempt isolates
a summable-scale criterion and constructs a nonstationary current
showing why density and rectifiability alone cannot replace it.

\paragraph{What monotonicity actually controls.}
Translate the point to the origin and remain inside a ball disjoint
from the boundary. The varifold associated with a minimizing current
is stationary. Write $\Theta(r)=r^{-m}\|T\|(B_r)$.
With this normalization the Euclidean monotonicity identity is
\begin{equation}\label{a418:mono}
 \Theta(R)-\Theta(r)=
 \int_{B_R\setminus B_r}
 \frac{|x^\perp|^2}{|x|^{m+2}}\,d\|T\|(x).
\end{equation}
The identity holds at admissible radii and extends by limits.
Here $x^\perp$ is normal to the approximate tangent plane.
Consequently $\Theta(0)$ exists and the radial tilt energy
is finite near zero.

For $r_j=2^{-j}r_0$, put
$e_j=\Theta(r_j)-\Theta(r_{j+1})\geq0$. Then
\[
                  \sum_{j\geq0}e_j=\Theta(r_0)-\Theta(0)<\infty.
\]
This scalar energy statement does not imply that the cones
selected at successive scales have summable total change.
Even an estimate of angular change by $\sqrt{e_j}$, if
justified in the relevant setting, would require summability
of those square roots. For example $e_j=(j+1)^{-2}$ is
summable and $\sum_j\sqrt{e_j}$ diverges. Rotations by
$1/(j+1)$ can keep turning while their squared angular
increments have finite sum. Applying Cauchy--Schwarz
over infinitely many scales would introduce the square
root of an infinite number of scales, not a finite bound.

\paragraph{A precise Cauchy criterion for full currents.}
On locally mass-bounded currents choose a metric $d$ for
local weak convergence using a countable dense family of
compactly supported smooth test forms. A weighted sum
of $\min\{1,|(S-T)(\omega_k)|\}$ suffices, with the family
dense on every compact set. Uniform local mass bounds
make convergence on the family equivalent to convergence
against every test form.

Suppose
\begin{equation}\label{a418:oscillation}
 \sup_{r,s\in[r_{j+1},r_j]}d(T_{0,r},T_{0,s})\leq\eta_j,
 \qquad \sum_j\eta_j<\infty.
\end{equation}
Then every blow-up sequence has the same limit. Indeed the
dyadic sequence is Cauchy, since successive distances are
at most $\eta_j$. Local compactness of minimizing integral
currents supplies a convergent subsequence; hence the
whole dyadic sequence converges to its limit $C$.
For an arbitrary intermediate radius,
\[
 d(T_{0,r},C)\leq\eta_j+d(T_{0,r_j},C)\longrightarrow0
 \quad\text{if }r\in[r_{j+1},r_j].
\]
This proves convergence through all radii, not merely
along a chosen discrete sequence.

An actual comparison $\eta_j\leq A\sqrt{e_j}$, together
with $e_j\leq B2^{-\beta j}$ for $\beta>0$, would be a
sufficient package: the resulting geometric series converges.
Neither that comparison in this metric nor such power
decay is established here for every singular point.
They are additional hypotheses. Monotonicity alone
supplies only the sum of $e_j$.

\paragraph{A complete regular-point calculation.}
Suppose near zero the current is
$Q\lbrack\!\lbrack\operatorname{graph}u\rbrack\!\rbrack$ over an oriented
$m$-plane, with $u$ of class $C^1$, $u(0)=0$ and $Du(0)=0$.
At scale $r$ the graphing function is $u_r(y)=u(ry)/r$.
On fixed compact subsets,
\[
 \sup|u_r|\longrightarrow0,\qquad
 \sup|Du_r|=\sup|Du(ry)|\longrightarrow0.
\]
Parametrized integrals of compactly supported test forms
converge to those on the plane, retaining its orientation
and the integer coefficient $Q$. The tangent current
is uniquely $Q\lbrack\!\lbrack P\rbrack\!\rbrack$. If
$|Du(x)|\leq C|x|^\alpha$, the same calculation gives
a power rate for fixed test-form norms and exhibits
the summable-scale mechanism explicitly.

This argument assumes a genuine smooth graph near
the point. A tangent plane with multiplicity $Q$
at a singular point does not automatically provide
such a graph. Replacing that hypothesis by its
desired conclusion would conceal the regularity gap.

\paragraph{A slowly rotating current as a failure test.}
For $0<r<r_0<1$, let
\[
 \phi(r)=\log\log(e/r),\qquad
 \gamma(s)=s(\cos\phi(|s|),\sin\phi(|s|))
 \quad(s\ne0),\qquad\gamma(0)=0.
\]
The parametrization is Lipschitz, since
\[
 |\gamma'(s)|=\sqrt{1+\log(e/|s|)^{-2}}
\]
is bounded. It defines an integral one-current with no
boundary at zero. Because $|\gamma(s)|=|s|$, its density is
\[
 \frac{\|\gamma_\#\lbrack\!\lbrack(-r_0,r_0)\rbrack\!\rbrack\|(B_r)}{2r}
 =\frac1r\int_0^r\sqrt{1+\log(e/s)^{-2}}\,ds
                                      \longrightarrow1.
\]
The scalar density therefore agrees with that of a
multiplicity-one line.

Choose $r_j\downarrow0$ for which $\phi(r_j)$ tends
modulo $2\pi$ to $\theta$. For fixed $s\ne0$,
\[
 \phi(r_j|s|)-\phi(r_j)
 =\log\left(1+\frac{\log(1/|s|)}{\log(e/r_j)}\right)
                                      \longrightarrow0.
\]
The rescaled parametrizations and derivatives converge
away from zero to $s\mapsto s(\cos\theta,\sin\theta)$.
Their uniform derivative bound permits dominated
convergence in test-form integrals across zero.
A test form supported in $B_R$ only samples
parameters $|s|\leq R$, so no parameter tail is
lost. The tangent current is the oriented line
of angle $\theta$, and different angle subsequences
give different tangent currents.

This is deliberately not a minimizing example.
The tangent angle changes along every punctured
branch, giving nonzero curvature and a nonzero
first variation for a suitable compactly supported
normal variation away from zero. Thus it is not
even stationary. Its product with $\mathbb R^{m-1}$,
with one extra zero ambient coordinate if necessary,
has the same tangent nonuniqueness in the dimensional
range of the target, but still fails stationarity.
It does not satisfy the full hypotheses of
\eqref{a418:mono}. The example tests a shortcut
and is not a counterexample to the conjecture.

\paragraph{Density, support and exceptional points are different issues.}
All tangent cones have the same density, but distinct
oriented cones may share one density. Even equality
of support would not alone identify orientations
and multiplicities. A selection argument must
compare currents or quantities controlling their
test-form evaluations, as in
\eqref{a418:oscillation}; density is not such a metric.

Likewise uniqueness at almost every point of a
rectifiable singular stratum does not settle
every interior point. A full-measure subset
may leave exceptional points at which the
available decay theorem says nothing. The
universal quantifier in the target cannot
discard those points. This is why the strong
stratified results described in [S1] do not
by themselves close the final step here.

\paragraph{Verification, first gap, and outcome.}
The Cauchy argument includes every intermediate radius,
and the rotating example is checked as an oriented
current rather than as a picture of supports.
Its speed is bounded while its total turning
diverges, matching the distinction between
squared increments and their sum. At regular
points the multiplicity is preserved under
every rescaling.

The first missing general input is a sufficiently
strong decay or cone-selection estimate at each
singular point. A concrete continuation would
seek the within-annulus control in
\eqref{a418:oscillation}, together with a
summability mechanism stronger than the scalar
monotonicity budget.
\textbf{UNRESOLVED.} Regular graph uniqueness
and the conditional scale criterion are
proved. The explicit nonunique example
violates the variational hypothesis, and
no minimizing counterexample or general
uniqueness theorem is claimed.

\subsection{Sources}
\begin{itemize}
\item[S1]Skorobogatova, Singularities of area-minimizing integral currents: recent progress and open directions, Seminaire Laurent Schwartz EDP et applications (2025).
\url{https://proceedings.centre-mersenne.org/item/10.5802/slsedp.175.pdf}.
\textit{Formulation source.}
\end{itemize}
\end{document}
